\documentclass[12pt,DIV=14, version=first, BCOR=10mm,a4paper,twoside,parskip=half-,headsepline,headinclude]{scrartcl}
% Grundgröße 12pt, zweiseitig
% Standardpakete
\usepackage[utf8]{inputenc}
\usepackage[T1]{fontenc}
\usepackage{lmodern}
% deutsche Silbentrennung
\usepackage[ngerman]{babel}
% schöne Tabellen
\usepackage{booktabs}
% etwas Mathematik
\usepackage{amsmath,amssymb}
% Grafiken einbinden
\usepackage{graphicx}
% Kopf- und Fußzeilen
\usepackage[headsepline,footsepline,automark]{scrlayer-scrpage}
\pagestyle{scrheadings}
% Links anklickbar
\usepackage{url,hyperref,color}
\definecolor{MidnightBlue}{cmyk}{0.98,0.13,0,0.43}
\hypersetup{colorlinks,breaklinks=true, pdffitwindow=true,pdfpagelayout=SinglePage,
            linkcolor=blue,pdfauthor={}, pdfdisplaydoctitle=true,
            pdfsubject={},   % to be defined later
            bookmarksnumbered,citecolor=MidnightBlue,
            urlcolor=MidnightBlue}
\raggedbottom
\renewcommand{\topfraction}{1}
\renewcommand{\bottomfraction}{1}


% zum Ausprobieren
\usepackage{blindtext}


% jetzt gehts los
\begin{document}% hier gehts los
  \thispagestyle{empty} % Titelseite
\includegraphics[width=0.2\textwidth]{Wortmarke_WI_schwarz}

   {  ~ \sffamily
  \vfill
  {\Huge\bfseries Ruby als Lernprogrammiersprache}
  \bigskip

  {\Large
  Nicolai Böttger und Jon-Steven Streller \\[2ex]
 Seminar-Arbeit im Studiengang "`Angewandte Informatik"'
 \\[5ex]
   \today }
}
 \vfill

  ~ \hfill
  \includegraphics[height=0.3\paperheight]{H_WI_Pantone1665}

\vspace*{-3cm}

  \newpage \thispagestyle{empty}
 \begin{tabular}{ll}
{\bfseries\sffamily Autor 1:} &  Nicolai Böttger \\
            & 1476431 \\
            & \href{mailto:nicolai.boettger@stud.hs-hannover.de}{nicolai.boettger@stud.hs-hannover.de} \\
            & Verfasste Seiten/Abschnitte: ...
            \\[5ex]
{\bfseries\sffamily Autor 2:} & Jon-Steven Streller \\
            & 1475759 \\
            & \href{mailto:steven.streller@stud.hs-hannover.de}{steven.streller@stud.hs-hannover.de} \\
           & Verfasste Seiten/Abschnitte: ... \\[5ex]
 {\bfseries\sffamily Prüfer:} &Prof. Dr. Dennis Allerkamp \\
          & Abteilung Informatik, Fakultät IV \\
         & Hochschule Hannover \\
        & \href{mailto:dennis.allerkamp@hs-hannover.de}{dennis.allerkamp@hs-hannover.de}
\end{tabular}

\vfill

\begin{center} \sffamily\bfseries Selbständigkeitserklärung \end{center}
% fett und zentriert in der minipage

Mit der Abgabe der Ausarbeitung erklären wir, dass wir die eingereichte Seminar-Arbeit
selbständig und ohne fremde Hilfe verfasst, andere als die von uns angegebenen Quellen
und Hilfsmittel nicht benutzt und die den benutzten Werken wörtlich oder
inhaltlich entnommenen Stellen als solche kenntlich gemacht haben.
\vspace*{7ex}

Hannover, den \today \hfill

\pagebreak
\section{Abstract}
\begin{flushleft}
Dieser Abstrakt fasst den Inhalt und die Ergebnisse dieser Arbeit kurz zusammen. Diese Arbeit beschäfftigt sich mit der Frage, ob sich die Programmiersprache \textit{\glqq Ruby\grqq} als erste Programmiersprache für den Einsatz in einer Schule, oder ähnlichem, eignet. Die aufgestellten Kriterien - Einstiegsfreundlichkeit, Skalierbarkeit, Verständlichkeit, Dokumentation, Verbreitung der Sprache und Anforderungen an die Schule - welche zum Beantworten dieser Frage dienen sollten, wurden für \textit{\glqq Ruby\grqq} überwiegend zutreffend bewertet, weswegen diese Arbeit zu dem Ergebnis kommt, dass sich \textit{\glqq Ruby\grqq} grundsätzlich schon für diese Aufgabe eignet. Trotzdem, ist zu beachten, dass dies nicht heißen soll, dass Ruby die perfekte erste Programmiersprache wäre.
\end{flushleft}
  \pdfbookmark[0]{Inhalt}{contents}
  \tableofcontents  % Inhaltsverzeichnis

\pagebreak

%\input{abkuerz.tex}      % Einbinden von Tex-Files
%\input{einfuehrung.tex}


% durch eigenen Text ersetzen

\section{Einführung}
\begin{flushleft}
Die vorliegende Arbeit entstand im Rahmen des Seminars \textit{\glqq Programmiersprachen für den Einstieg\grqq}. Ziel des Seminars war es mehrere Programmiersprachen auf ihre Tauglichkeit im Anwendungsbereich einer pädagogischen Lehrveranstaltung zu analysieren. Hierbei wurde jede Programmiersprache einzeln von verschiedenen Studierenden in Zweiergruppen vorgestellt und die jeweiligen Ergebnisse Ihrer Recherche präsentiert. Nachfolgend entstand ein Austausch zwischen den Studierenden über positive und auch negative Auffälligkeiten in der Präsentation.
\end{flushleft}

\subsection{Vorgehensweise}
\begin{flushleft}
Um eine sinnvolle Bewertung der Sprache in Bezug zum Thema des Seminars zu erreichen haben wir zuerst nach Studien gesucht, welche Kriterien definieren, nach denen man eine Programmiersprache bewerten kann. Dazu haben wir zuerst nach wissenschaftlichen Arbeiten gesucht, welche sich schon einmal mit diesem Thema befasst haben. Da wir dadurch aber nicht ausreichend fündig wurden, haben wir danach nach Arbeiten gesucht, die sich mit ähnlichen Fragestellungen befasst haben und Kriterien dieser auf unser Thema übertragen. Danach haben wir angefangen uns selbst in unsere vorzustellende Sprache \textit{\glqq Ruby\grqq} einzuarbeiten. Nachdem wir ein grundsätzliches Verständnis der Sprache bekommen haben, haben wir angefangen \textit{\glqq Ruby\grqq} in Bezug zu den vorher gewählten Kriterien zu bewerten.
\end{flushleft}

\subsection{Aufbau der Arbeit}
\begin{flushleft}
Die Arbeit beginnt damit, dass zuerst die Anforderungen an die Schule und die Schüler in Bezug zum Lernen einer Programmiersprache dargestellt werden. Als nächstes wird kurz die Sprache \textit{\glqq Ruby\grqq} an und für sich vorgestellt und etwas auf ihre interne Arbeitsweise eingegangen. Danach wird die Methodik vorgestellt, nach der versucht wurde, die Fragestellung der Arbeit zu beantworten. Im Hauptteil werden dann die in Verbindung mit der Methodik entstandenen Kriterien definiert und die Programmiersprach \textit{\glqq Ruby\grqq} anhand jener bewertet. Anschließend daran folgt im Schlussteil zum einen noch eine zusammenfassende Bewertung der Sprache im Bezug zum Thema und zum anderen einige Schlussbemerkungen, die es beim Interpretieren der Ergebnisse dieser Arbeit und des durchgeführten Seminars an sich zu beachten gibt.
\end{flushleft}

\subsection{Anforderungen}
\begin{flushleft}
Die nun genannten Anforderungen sind durch die Recherche und Interpretierung  der Forschungsfrage entstanden. Grundlegend sollte es bei der Zielgruppe um mindestens eine achte oder neunte Klasse handeln, da viele Programmierausdrücke in der Programmiersprache sehr ähnlich oder gar gleich zur Mathematik sind. Eine gewisse Reife ist Vorauszusetzen, um gewährleisten zu können, dass der/die Schüler/-in sich die Auswirkung aufs zukünftige Berufsleben in Hinsicht der erworbenen Qualifizierung bewusst ist um somit auch ein Interesse an der Thematik zu entwickeln. Außerdem sollte das Prinzip von Logik zumindest vorhanden sein. Vorkenntnisse im Bereich der Entwicklung von Anwendung wäre wünschenswert aber keine Anforderung an den Schüler. Allerdings sollte die Bedienung eines Computers keine Probleme darstellen und grundlegende Steuerung wie zum Beispiel mit der Maus sollte bereits beherrscht werden.
\end{flushleft}

\section{Hauptteil}

\subsection{Was ist Ruby?} \label{Was ist Ruby?}
\begin{flushleft}
Im Jahr 1995 veröffentlichte \textit{Yukihiro Matsumoto} die erste Version (Version 0.95) von \textit{\glqq Ruby\grqq}. \textit{Yukihiro Matsumoto} lies sich während der Entwicklung von \textit{\glqq Ruby\grqq} von mehreren Programmiersprachen wie z.B. Perl, Smalltalk, Lisp, Ada inspirieren. \textit{Ruby} besitzt eine tief integrierte
Objektorientierung\footnote{\label{foot:1} Ein System besteht in der Objektorientierung ausschließlich aus Objekten, die miteinander über Nachrichten kommunizieren. Jedes Objekt verfügt über Eigenschaften und Methoden. Die Eigenschaften beschreiben dabei über ihre Werte den Zustand eines Objektes, die Methoden die möglichen Handlungen eines Objektes. vgl. [WL01]} (kurz OO). Ruby wird ungefähr in einem Intervall von 3 Monaten mit Updates versorgt. Im Moment (17. April 2019) befindet sich Ruby in Version 2.6.3 Die Version 2.7.0 steht schon in einer Vorabversion zur freien Verfügung. Die Programmiersprache Ruby ist Plattformunabhängig und somit auf kein bestimmtes System beschränkt. Trotz vollständiger Objektorientierung können aber auch Programmierparadigmen wie prozedurale oder funktionale Programmierung verwendet werden.
\end{flushleft}

\subsection{Arbeitsweise von Ruby}
\begin{flushleft}
An dieser Stelle wird kurz auf die interne Arbeitsweise von Ruby eingegangen. Die Programmiersprache Ruby basiert auf einen \textit{Interpreter}. Um ein besseres Verständnis zu vermitteln, wird hier auf die Unterscheidung zwischen Kompilierer (engl. Compiler) und Interpreter eingegangen. Der \textit{Interpreter} arbeitet Zeile für Zeile den Quellcode aus und stoppt nur wenn ein Fehler in einem Ausdruck vorgefunden wurde, während bei dem Kompilierer direkt der ganze Quellcode analysiert wird und somit direkt alle Fehler aufgelistet werden. Wie bereits vermuten lässt, ist der Kompilierprozess deutlich anfordernder für die Ressourcen als der \textit{Interpreter}. Für den \textit{Interpreter} spricht, dass die Analyse sehr schnell ausgeführt wird und weniger Leistung erfordert. Während bei dem Kompilierer weitere Bibliotheken in den Quellcode eingefügt werden ist dies bei dem \textit{Interpreter} nicht nötig. Dadurch ist der \textit{Interpreter} in Hinsicht auf die Speichernutzung etwas effektiver.
\end{flushleft}

\subsection{Methodik}
\begin{flushleft}
Um die Programmiersprache Ruby in Bezug zur Fragestellung, ob sie sich als Programmiersprache für den Einstieg eignet bewerten zu können, wurde als Methodik ein Kriterienkatalog erstellt, welcher die Sprache in hinsichtlich verschiedener Bereiche untersucht. Dieser Katalog besteht aus den Kriterien Einstiegsfreundlichkeit, Skalierbarkeit, Verständlichkeit, Dokumentation, Verbreitung der Sprache und Anforderung an Lehranstalt. Eine besondere Gewichtung wurde hierbei auf die Kriterien Skalierbarkeit und Verständlichkeit gelegt, da diese die grundlegenden Aspekte zum Verstehen der Sprache sind. Nach einer kurzen Erläuterung der Kriterien, werden diese in Bezug zur Programmiersprache Ruby gesetzt und es wird erörtert, inwiefern das Kriterium zutrifft (trifft voll zu, trifft teilweise zu, trifft nicht zu)
\end{flushleft}

\subsubsection{Einstiegsfreundlichkeit}

\begin{flushleft}
Das Kriterium der \textit{Einstiegsfreundlichkeit} fast alles zusammen, was benötigt wird um mit dem Programmieren in der Sprache anzufangen. So sollte die Installation von möglichen Compilern oder Interpretern nicht zu aufwändig sein und ohne große Fachkentnisse vorgenommen werden können. Auch wäre es von Vorteil, wenn die Installationsbestandteile nicht allzu viel Speicher einnehmen. Ebenfalls wird hier berücksichtigt, ob beispielsweise mehrere Dateien angelegt, oder andere Vorgänge wie ein mögliches Linken von Dateien geschehen muss, ehe man ein geschriebenes Programm ausführen kann.

Die technischen Vorraussetzungen, die benötigt werden, um mit der Programmiersprache \textit{\glqq Ruby\grqq} arbeiten zu können beschränken sich ausschließlich auf einen \textit{Ruby-Interpreter}, welchen man auf der offiziellen Website von \textit{\glqq Ruby\grqq} herunterladen kann. Auch die anschließende Installation läuft unter dem Betriebssystem \textit{Microsoft Windows} so ab, wie man es von anderen heruntergeladenen Programmen gewohnt ist. Die Größe der Standardedition dieses Interpreters beläuft sich auf ca. 60MB und nimmt somit vergleichsweise wenig Speicher auf der Festplatte ein (Java SE Development Kit 12.0.1 ~160MB).
Theoretisch könnte man in \textit{\glqq Ruby\grqq} ein komplettes Programm in nur einer Datei schreiben, welches man ohne weitere Zwischenschritte von einem Terminal aufrufen könnte. Zusätzlich wird mit dem Download des Interpreters das Programm \textit{IRB} (Interactive Ruby) installiert, mit welchem man unmittelbares Feedback auf eingegebenen Ruby-Code erhält.

Dementsprechend trifft das gewählte Kriterium der \textit{Einstiegsfreundlichkeit} voll zu.
\end{flushleft}

\subsubsection{Skalierbarkeit} \label{Skalierbarkeit}
\begin{flushleft}
Unter Skalierbarkeit wird alles bewertet, was mit den spezifischen Eigenschaften und Besonderheiten der Sprache zu tun hat. Wie hoch ist die Komplexität um eine Ausgabe zu erzeugen? Wie sehr bindet mich die Programmiersprache an bestimmte Syntax und Semantik? Ein weiterer Punkt der nicht zu vernachlässigen ist, ist die Plattformunabhängigkeit der Sprache. Hierzu wird im nachfolgenden Abschnitt näher drauf eingegangen.

Durchaus positiv sind die herrschenden Konventionen in \textit{\glqq Ruby\grqq}. Hierbei lässt die Sprache selbst sehr viel Spielraum und Akzeptanz bei der Syntax und Semantik. Dies äußert sich sehr stark bei der Verwendung von Semikolon - während es bei vielen etablierten Programmiersprachen strikt und nur eine Richtlinie gibt, wird bei \textit{\glqq Ruby\grqq} ein ganz anderer Weg eingeschlagen. Es ist möglich, aber nicht nötig Semikolons zu setzen. Es ist auch nicht nötig sich für einen Stil zu entscheiden, so kann jeder seinen eigenen Programmierstil in dem Lehrprozess entwickeln. Wie in Abschnitt \ref{Verständlichkeit} bereits erwähnt, ist das Konzept der Objektorientierung nicht zwingend zum anwenden von \textit{\glqq Ruby\grqq} erforderlich. Dadurch ist es von \textit{\glqq Ruby\grqq} selbst gegeben, bereits mit nur \textbf{einer} Zeile Code eine Ausgabe auf der Kommandozeile\footnote{\label{foot:2} {Kommandozeile: Hierbei handelt es sich um eine grundlegende Anwendung die meist durch Texteingaben gesteuert wird.}} zu wiedergeben. Dieser Effekt wird erzielt weil Klassen in \textit{\glqq Ruby\grqq} nicht erforderlich sind und somit sehr viel Quellcode bei einfachen Anwendungen eingespart werden kann.

In den meisten Fällen haben Programmiersprachen nur einen oder wenig Anwendungsgebiete. \textit{\glqq Ruby\grqq} sticht hierbei etwas heraus, denn die Vielseitigkeit der Programmiersprache ist enorm. \textit{\glqq Ruby\grqq} besitzt die Möglichkeit als Bash-Anwendung\footnote{\label{foot:3} {Bash-Anwendung: automatisiertes Programmieren auf der Kommandozeile}} und separat in der Webentwicklung zu fungieren. Dies ist ein exemplarisches Beispiel wie stark der Kontrast beziehungsweise die Einsatzmöglichkeiten von \textit{\glqq Ruby\grqq} sind.

Das erweitern von Bibliotheken beziehungsweise Paketen gestaltet sich in \textit{\glqq Ruby\grqq} denkbar einfach. Hier punktet \textit{\glqq Ruby\grqq} mit einem zentralen Paketverwaltungssystem (Link: \href {https://rubygems.org/} {Ruby Gems})  in dem zusätzlich für jedes Paket eine Installationsanleitung sowie eine Dokumentation und Quellcodebeispiele zum verdeutlichen der Arbeitsweise des Paketes beinhaltet. Hierbei muss jedoch beachtet werden, dass jeder die Möglichkeit hat Pakete selbst zu veröffentlichen und somit nicht verifiziert werden kann, dass jedes Paket fehlerfrei beziehungsweise auf einem gewissen Niveau entwickelt wurde. Allerdings ist ein hoher Download-Aufruf ein Indikator dass das Paket einem gewissen Standard entspricht. Hierbei kann \textbf{nur} von einer Korrelation zwischen Funktionalität des Quellcodes und Download-Anzahl ausgegangen werden!

Negativ anzumerken ist, dass Programmiersprachen die einen Interpreter verwenden (so wie \textit{\glqq Ruby\grqq}) in der Ausführung langsamer sind als Programmiersprachen mit einem Kompilierer (engl. Compiler).

Trotz der Tatsache dass die Sprache eine hohe Toleranz in Bezug auf \textit{Skalierbarkeit} sich hierbei aufweist, trifft das Kriterium wegen der eventuellen schlechten Qualität der zu verfügbar gestellten Pakete und aufgrund der Nutzung eines \textit{Interpreters} in Zusammenhang des Geschwindigkeitsdefizit nur teilweise zu.

\end{flushleft}

\subsubsection{Verständlichkeit} \label{Verständlichkeit}
\begin{flushleft}
Bei der \textit{Verständlichkeit} geht es vor allem darum, dass der geschriebene Quellcode möglichst simpel und verständlich wirkt. Ein abstrakter Quellcode wäre nicht fördernd um grundlegende Konzepte des Programmieren zu vermitteln. Außerdem ist ein schnelles und klares Feedback des Interpreters essentiell, um angenehm und ohne großes Suchen nach Fehlern programmieren zu können.

Die Programmiersprache \textit{\glqq Ruby\grqq} besitzt eine sehr hohe \textit{Verständlichkeit}. Auch wenn \textit{\glqq Ruby\grqq} eine tief integrierte Objektorientierung (nachstehend \glqq OO\grqq) \, besitzt, ist es nicht zwingend erforderlich das Konzept der eigentlichen OO zu verstehen. In \textit{\glqq Ruby\grqq} ist es auch nicht nötig einen expliziten Einstiegspunkt zu definieren, wie es zum Beispiel in Java mit der klassischen Main-Methode der Fall ist. Was das erlernen der Programmiersprache \textit{\glqq Ruby\grqq} auch deutlich vereinfacht, ist dass \textit{\glqq Ruby\grqq} an dem Englischen angelehnt ist und der gesprochenen Sprache sehr ähnelt. Dies äußert sich öfter in Quellcodebeispielen, bei denen man einzelne Abschnitte wie gesprochene Sätze ablesen kann. Der \textit{Interpreter} hilft aktiv bei der Fehlersuche im Quellcode und somit ist der Grad der Frustration bei nichtfunktionierenden Code geringer als bei Programmiersprachen die auf einen \textit{Kompilieren} (engl. Compiler) setzen.

Dementsprechend trifft das gewählte Kriterium der \textit{Verständlichkeit} voll zu.
\end{flushleft}

\subsubsection{Dokumentation}
\begin{flushleft}
Eine aktuelle und feinsäuberlich ausgearbeitete \textit{Dokumentation} der Programmiersprache ist überlebenswichtig für jene. Beachtet werden muss auch zum Beispiel wie viel Lehr/-Sachbücher über die Programmiersprache publiziert wurden. Ergänzend hierzu ist es immer sehr positiv, wenn die Programmiersprache über viele Entwickler/-innen verfügt, da so ein großer Support bei technischen Fragen eher gegeben ist. Damit Schüler auch von Zuhause aus programmieren können, wäre es hilfreich, wenn spezifische Webseiten mit Tutorials oder Lernvideos zum Erlernen dieser Programmiersprache existieren würden.

Die \textit{\glqq Ruby-Doc\grqq} (Abkürzung von \glqq Ruby-Dokumentation\grqq) (\href {https://ruby-doc.org/}{Hier} einsehbar) besitzt eine sehr ordentliche Strukturierung und weist ein sehr klares und verständliches Layout auf. Gerade für Anfänger ist es wichtig eine sauber ausgearbeitete Dokumentation vorzufinden, damit der Lehrprozess hier nicht durch Verständnisprobleme wegen mangelnder Erläuterung unterbrochen wird. Hierbei finden sich in der \textit{\glqq Ruby-Doc\grqq} Anwendungsbeispiele die zur Unterstürzung des Verständnis beitragen. \\
Hinsichtlich Literatur, also Lehr- und Sachbüchern lässt die Dokumentation der Sprache \textit{\glqq Ruby\grqq} zu wünschen übrig. Es gibt zwar einige Bücher, diese sind aber in der Regel schon etwas älter und dementsprechend nicht mehr kompatibel mit dem heutigen Stand der Sprache. Besonders schwierig ist es gute Literaturbücher in der deutschen Sprache zu finden (\href {https://wiki.ruby-portal.de/Literatur.html} {Hier} einsehbar). Im Gegensatz dazu gibt es aber einige gute Videotutorials, entweder kostenlos auf Youtube oder kostenpflichtig - teilweise mit Übungsaufgaben - auf Plattformen wie {\glqq Udemy\grqq} um sich das Programmieren mit \textit{\glqq Ruby\grqq} selbst beizubringen.
Trotz der Tatsache, dass man die Grundlagen der Programmiersprache so auch gut über das Internet lernen kann, trifft das Kriterium einer ausgereiften \textit{Dokumentation}, aufgrund fehlender schriftlicher Literatur nur teilweise zu.
\end{flushleft}

\subsubsection{Verbreitung der Sprache}
\begin{flushleft}
Unbeliebte Programmiersprachen sind in der Regel auch nicht so stark in der Arbeitswelt gefordert, weshalb es sehr ernüchternd sein kann, eine Programmiersprache pädagogisch vermittelt zu bekommen die keine reelle Perspektive in der Zukunft besitzt. Zusätzlich ist eine Würdigung beziehungsweise eine Akzeptanz in der \glqq Programmiergesellschaft\grqq \, wünschenswert, da es durchaus vorkommt, dass Programmiersprachen sich eher schwer durchsetzen, wenn diese nicht von der o.g Zielgruppe akzeptiert werden. Desweiteren ist eine aktive Weiterentwicklung der Programmiersprache eine elementare Voraussetzung um gegen andere aktive und etablierte Sprachen zu konkurrieren.
Anfänglich lässt sich sagen, dass \textit{\glqq Ruby\grqq}, wie auch schon im Kriterium der Skalierbarkeit erklärt, grundsätzlich vielseitig einsetzbar ist. Auch lässt die Statistik von {\glqq RedMonk\grqq} ([BP 01]), die die größten Programmiersprachen nach Beliebtheit auf {\glqq Stack Overflow\grqq} (Anzahl Tags) und GitHub (Anzahl Projekte), in der \textit{\glqq Ruby\grqq} in beiden Bereichen gut abschneidet, vermuten, dass die Sprache zumindest eine große Community besitzt. Trotz dieser vielseitigen Einsatzmöglichkeit sieht es in der Praxis aber doch meist so aus, dass \textit{\glqq Ruby\grqq} hauptsächlich in Verbindung mit dem Framework {\glqq Ruby on Rails\grqq} zur Entwicklung von Webanwendungen benutzt wird. Trotzdem, nimmt sie im Vergleich zu Sprachen wie {\glqq PHP\grqq} oder {\glqq Asp.net\grqq} nur einen kleinen Teil der Webseiten ein ([RW 01]) . Auf der anderen Seiten wird die Sprache aber in relativ regelmäßigen Abständen seit 2002 weiterentwickelt ([RR 01]). Zusammenfassed lässt sich also festhalten, dass \textit{\glqq Ruby\grqq} zwar eine aktive Community hat, in der Praxis aber nicht allzu oft verwendet wird.

Dementsprechend trifft das gewählte Kriterium der \textit{Verbreitung der Sprache} teilweise zu.
\end{flushleft}

\subsubsection{Anforderung an Lehranstalt}
\begin{flushleft}
Selbstverständlich sollten die Lehrmittel beziehungsweise Lizenzkosten im Optimalfall keine Kosten verursachen, um auch Schulen mit wenig Budget eine Perspektive bieten zu können. Ideal wäre ein/e Lehrbeauftragter/-in die bereits grundlegende Kenntnisse in der Programmiersprache beherrscht. Somit könnten Ausbildungsmaßnahmen für den/die Lehrbeauftragte/-n eingespart werden und die eingesparten Kosten in die Lehrveranstaltung investiert werden.

\textit{\glqq Ruby\grqq} ist eine frei verfügbare Sprache und ist mit keinerlei Lizenzkosten an den/die Entwickler/-in geknüpft. Durchaus positiv ist auch die simple Verwaltung beziehungsweise Wartung von \textit{\glqq Ruby\grqq}. Hierbei wird kein externer Server benötigt, wofür weitere Fachkenntnisse erforderlich wären. Somit sind die Kosten und der Aufwand für die Benutzung von \textit{\glqq Ruby\grqq} niedrig, da kein weiteres Equipment außer der Computer erforderlich ist. Da \textit{\glqq Ruby\grqq} wie bereits in \ref{Was ist Ruby?} erwähnt auf größere Programmiersprachen basiert beziehungsweise von diesen inspiriert wurde, besteht eine größere Chance eine Lehrkraft zu finden die bereits die Grundlagen beherrscht. Somit kann man davon ausgehen, dass die Lehrkraft einen schnellen Einstieg nach kurzer Einarbeitungsphase in die Besonderheiten von \textit{\glqq Ruby\grqq} absolvieren und somit die Schüler unterrichten kann. Eine permanente Internetverbindung wäre von Vorteil, da es nicht auf einfachem Wege ist externe Pakete lokal über ein Wechselmedium oder ähnliches zu installieren. Hierbei muss wie schon bereits in Abschnitt \ref{Skalierbarkeit} erwähnt, ist die Paketinstallation nur über das zentrale Paketverwaltungssystem möglich, und somit müsste jeder Computer der von einem/einer Schüler/-in genutzt wird, zumindest während der Installation des Paketes über eine funktionierende Internetverbindung verfügen.

Auch wenn man den Punkt mit der Internetverbindung durchaus bemängeln kann, ist dies kein grundlegendes Problem welches die Lernzielerfolge beeinträchtigen würde, da gerade bei Anfänger/-innen teilweise oder auch komplett auf das Paketverwaltungssystem verzichtet werden kann um grundlegende Programmiertechniken zu lehren. Dementsprechend trifft das gewählte Kriterium \textit{Anforderung an Lehranstalt} trotz des kleines Mängel voll zu.
\end{flushleft}

\section{Bewertung}
\begin{flushleft}
In diesem Abschnitt wird eine abschließende Bewertung, zu der Frage, ob sich \textit{\glqq Ruby\grqq} als erste Programmiersprache für den Einsatz im Unterricht eignet zusammenfassed dargestellt. Die jeweiligen Bewertungen der einzelnen Kriterien führen zu einer Gesamtbewertung der Sprache von drei mal
{\glqq trifft voll zu\grqq} (Einstiegsfreundlichkeit, Verständlichkeit, Anforderung an Lehranstalt) und drei mal {\glqq trifft teilweise zu\grqq} (Skalierbarkeit, Dokumentation, Verbreitung der Sprache). Im Durchschnitt schneidet \textit{\glqq Ruby\grqq} also eher positiv ab.
Trotzdem ist dieses Ergebnis mit Vorsicht zu genießen, da es eigentlich notwendig ist, die Sprache an dieser Stelle mit anderen zu vergleichen, damit diese Bewertungen eine nachvollziehbare Aussagekraft bekommen.
\end{flushleft}

\section{Schlussbemerkungen}
\begin{flushleft}
Bei der eigentlichen Forschungsfrage hatten die Forscher in dem Fall die Studierenden freie Wahl bei der Auswahl der Kriterien und diese waren in der Realität auch öfter sehr unterschiedlich. In manchen Gruppierungen war zum Beispiel das Kriterium \textit{Kosten} eine relativ starke Gewichtung, während es in anderen Präsentationen gar keine Erwähnung fand. Der generelle Vergleich der verschiedenen Arbeiten gestaltet sich ohnehin schon schwer weshalb nur ein empirischer Vergleich in Frage kommt. Positiv ist allerdings eine breite und fast lückenlose Perspektive auf verschiedene Kriterien, als zum Beispiel eine Vorgabe des Prüfers. Auch wenn es in wissenschaftlichen Arbeiten eine hohe Objektivität Vorhanden sein soll, so ist es durch die verschiedenen Methoden der einzelnen Gruppen schwieriger eine \glqq objektive Skala\grqq zu definieren.
\end{flushleft}


\section{Literaturverzeichnis}
\begin{flushleft}
[WL01] Wirtschaftslexikon Gabler
\url {https://wirtschaftslexikon.gabler.de/definition/objektorientierung-43907}


[RW 01] Ruby in Webanwendungen
\url {https://i1.wp.com/www.pixelcrayons.com/blog/wp-content/uploads/2016/01/Most-popular-server-side-programming-languages.png?w=415{\&}ssl=1)}


[BP 01] Beliebtheit von Programmiersprachen
\url {https://redmonk.com/sogrady/2018/08/10/language-rankings-6-18/}


[RR 01] Ruby Releases\\
\url {https://www.ruby-lang.org/en/downloads/releases}
\end{flushleft}

% Literatur
\bibliographystyle{alpha}
%\bibliography{mybib}

\end{document}
